\subsection{INTEGRATING THE INHOMOGENEOUS LINEAR EQUATION}

Integrating the inhomogeneous linear equation

\[
{\frac {d \mathbf {x}}{d t}} = A (t). \mathbf {x} + \mathbf {b} (t)\tag{2}
\]

reduces to integrating the associated homogeneous equation

\[
{\frac {d \mathbf {x}}{d t}} = A (t). \mathbf {x}\tag{4}
\]

and evaluating a primitive. With the notation of th. 2 of IV, p.~179, let us put \(\mathbf{x} = C(t, t_0).\mathbf{z}\) , whence, from the second formula (6) of IV, p.~179, \(\mathbf{z} = C(t_0, t).\mathbf{x}\) ; if \(\mathbf{x}\) is an integral of (2) then \(\mathbf{z}\) is an integral of the equation \(\frac{d}{dt}\big(C(t, t_0).\mathbf{z}\big) = A(t)C(t, t_0).\mathbf{z} + \mathbf{b}(t)\) ; since the bilinear map

\[
(U, \mathbf {y}) \mapsto U. \mathbf {y}
\]

of \(\mathcal{L}(E) \times E\) into E is continuous (Gen.~Top., X, p.~297, prop. 6), z admits a derivative (except on a countable subset of J) and one has, by the formula for differentiating a bilinear function (I, p.~6, prop. 3)

\[
\frac {d}{d t} \big (C (t, t _ {0}). \mathbf {z} \big) = \frac {d C (t , t _ {0})}{d t}. \mathbf {z} + C (t, t _ {0}). \frac {d \mathbf {z}}{d t} = A (t) C (t, t _ {0}). \mathbf {z} + C (t, t _ {0}). \frac {d \mathbf {z}}{d t}
\]

(replacing \(dC(t, t_{0})/dt\) by \(A(t)C(t, t_{0})\) according to (5) (IV, p.~179)). The equation for z then reduces to \(C(t, t_{0}).d\mathbf{z}/dt = \mathbf{b}(t)\) , or again to

\[
\frac {d \mathbf {z}}{d t} = C (t _ {0}, t). \mathbf {b} (t)\tag{8}
\]

by the second formula (6) of IV, p.~179. Now the right-hand side of equation (8) is a regulated function on J, having been obtained by substituting the regulated functions U and y in the continuous bilinear function U.y (cf.~II, p.~55, cor. 2); equation (8) thus has one and only one integral taking the value \(x_{0}\) at \(t_{0}\) , given by the formula

\[
\mathbf {z} (t) = \mathbf {x} _ {0} + \int_ {t _ {0}} ^ {t} C (t _ {0}, s). \mathbf {b} (s) d s.\tag{9}
\]

Since one has \(C(t, t_{0})\) . \(\int_{t_{0}}^{t} C(t_{0}, s)\) . \(\mathbf{b}(s) \, ds = \int_{t_{0}}^{t} C(t, t_{0}) C(t_{0}, s)\) . \(\mathbf{b}(s) \, ds\) (II, p.~59, formula (9)), one obtains (taking account of the first formula (6) of IV, p.~179) the following result:

PROPOSITION 3. With the notation of th. 2 (IV, p.~179), for every point \((t_0, \mathbf{x}_0)\) of \(\mathbf{J} \times \mathbf{E}\) the integral of the linear equation (2) defined on \(\mathbf{J}\) and equal to \(\mathbf{x}_0\) at \(t_0\) is given by the formula

\[
\mathbf {u} (t) = C (t, t _ {0}). \mathbf {x} _ {0} + \int_ {t _ {0}} ^ {t} C (t, s). \mathbf {b} (s) d s.\tag{10}
\]

The method which leads to formula (10), which consists of taking the function z as a new unknown function, is often called the ``method of variation of parameters''.
% semantic-fragments: legacy-input-seam
\relax{}
